\chapter{Conclusiones y Trabajo Futuro}
\label{cap:conclusiones}

\chapterquote{Any fool can write code that a computer can understand. Good programmers write code that humans can understand.}{Martin Fowler}

\section{Conclusiones}

El objetivo principal de este trabajo era incorporar funcionalidades de realidad aumentada en tiempo real dentro de \textit{OBS Studio}, desarrollando un filtro nativo capaz de integrar contenido sintético directamente en el flujo de vídeo de una retransmisión. Ese objetivo se ha cumplido.

El resultado es un plugin funcional que extiende \textit{OBS Studio} con cuatro modos de operación diferenciados. En el modo de modelo 3D, el sistema carga geometría en múltiples formatos mediante \textbf{ASSIMP}, estima la pose del marcador ArUco en tiempo real usando el algoritmo PnP de \textbf{OpenCV} y renderiza el modelo anclado al marcador con la orientación y posición correctas, aplicando la cadena de transformaciones necesaria para convertir entre los convenios de coordenadas de OpenCV y OBS. En el modo cronómetro, ese mismo modelo actúa como reloj analógico de cuenta atrás cuyas manecillas se actualizan fotograma a fotograma y se sincronizan periódicamente con el tiempo oficial descargado de la API de \textbf{DOMjudge} \citep{DOMjudge}. En el modo Scoreboard, el sistema proyecta la clasificación general del concurso sobre el vídeo en tiempo real, con los datos de todos los equipos organizados en cuatro columnas. En el modo Team Info, la cámara detecta automáticamente qué equipo tiene en el campo de visión y muestra únicamente sus datos, combinando la detección de marcadores con un fichero JSON de mapeo que el operador prepara antes de la retransmisión.

Más allá de los modos de visualización, el trabajo ha resuelto varios problemas técnicos de envergadura. El sistema de coordenadas que separa OpenCV de OBS requirió una transformación explícita de convenios mediante una rotación de $180^\circ$ sobre el eje horizontal, con la consiguiente negación de las componentes $Y$ y $Z$ del vector de rotación ArUco. El suavizado del overlay mediante filtro de media exponencial, con tratamiento especial del salto de ángulo en la interpolación circular, elimina el temblor perceptible que producen el ruido de cámara y la compresión de vídeo. La arquitectura de dos hilos, con un hilo secundario dedicado exclusivamente a las peticiones HTTP y protegido por \texttt{CRITICAL\_SECTION}, garantiza que la red nunca bloquee el renderizado. El proceso de calibración de cámara, implementado como herramienta independiente en Python, genera el fichero de parámetros que el plugin carga al arrancar y que permite corregir la distorsión óptica de cada objetivo.

El sistema de conversión de formatos de vídeo cubre los siete formatos que OBS puede entregar (I420, NV12, YUY2, UYVY, BGRA, RGBA, Y800), haciendo que la detección funcione independientemente de la fuente de cámara conectada. La interfaz de usuario integrada en OBS permite al operador ajustar todos los parámetros en directo, sin reiniciar la escena, lo que hace el sistema manejable en un entorno de producción real.

En conjunto, el plugin demuestra que es posible añadir un sistema completo de realidad aumentada orientada a datos a \textit{OBS Studio} manteniendo la arquitectura de filtro nativo, sin latencia adicional por procesos externos y sin modificar el núcleo del software.

\section{Limitaciones}

El plugin ha sido desarrollado y probado en un entorno Windows con una cámara concreta. Esta circunstancia implica varias limitaciones que conviene señalar explícitamente. El sistema no ha sido validado en Linux ni en macOS, plataformas en las que OBS Studio también funciona pero cuya API de hilos y gestión de recursos difiere. La calidad de la detección depende directamente de las condiciones de iluminación y de la resolución de impresión de los marcadores, factores que no han sido evaluados de forma sistemática. La calibración de la cámara debe repetirse para cada objetivo diferente, lo que añade un paso de configuración que puede resultar poco práctico en entornos donde se cambia de cámara con frecuencia. Por último, el sistema de conectividad está diseñado exclusivamente para la API REST de \textbf{DOMjudge} \citep{DOMjudge}, por lo que no puede conectarse a otros sistemas de gestión de concursos de programación sin modificar el código.

\section{Trabajo Futuro}

Los resultados obtenidos abren varias líneas de continuación natural.

La más relevante desde el punto de vista técnico es la incorporación de \textbf{seguimiento sin marcadores} (\textit{markerless tracking}). Las técnicas actuales de SLAM visual \citep{MurArtal2015} permitirían anclar los \textit{overlays} a superficies del entorno real sin necesidad de imprimir y colocar códigos ArUco en las mesas, lo que haría el sistema más flexible y aplicable a escenarios donde no se puede controlar el espacio físico.

Una segunda línea es la \textbf{compatibilidad con otros sistemas de gestión de concursos}. Actualmente el cliente de red está adaptado a la API de DOMjudge; añadir soporte para plataformas como Codeforces, Kattis o el sistema ICPC CMS ampliaría el alcance del plugin a una fracción mucho mayor de la comunidad de concursos de programación.

En el plano del rendimiento, sería valioso realizar una \textbf{evaluación sistemática del impacto} del plugin sobre la tasa de fotogramas y el uso de CPU y GPU en diferentes configuraciones de hardware, estableciendo umbrales de uso recomendados y posibles optimizaciones para equipos de gama media.

Finalmente, añadir un \textbf{panel de control remoto} accesible desde una interfaz web o una aplicación móvil permitiría al director de realización gestionar los modos y la visibilidad de los overlays desde un dispositivo secundario, sin necesidad de tocar el equipo de retransmisión durante la emisión, lo que mejoraría significativamente el flujo de trabajo en producciones con más de una persona en el equipo técnico.
